
\documentclass{article}
\usepackage{colortbl}
\usepackage{makecell}
\usepackage{multirow}
\usepackage{supertabular}

\begin{document}

\newcounter{utterance}

\centering \large Interaction Transcript for game `hot\_air\_balloon', experiment `air\_balloon\_survival\_en\_reasoning off\_easy', episode 4 with deepseek{-}chat{-}v3.1{-}t1.0.
\vspace{24pt}

{ \footnotesize  \setcounter{utterance}{1}
\setlength{\tabcolsep}{0pt}
\begin{supertabular}{c@{$\;$}|p{.15\linewidth}@{}p{.15\linewidth}p{.15\linewidth}p{.15\linewidth}p{.15\linewidth}p{.15\linewidth}}
    \# & $\;$A & \multicolumn{4}{c}{Game Master} & $\;\:$B\\
    \hline

    \theutterance \stepcounter{utterance}  
    & & \multicolumn{4}{p{0.6\linewidth}}{
        \cellcolor[rgb]{0.9,0.9,0.9}{
            \makecell[{{p{\linewidth}}}]{
                \texttt{\tiny{[P1$\langle$GM]}}
                \texttt{Stai partecipando a un gioco di negoziazione collaborativa.} \\
\\ 
\texttt{Insieme a un altro partecipante, dovete accordarvi su un unico set di oggetti da conservare. Ognuno di voi ha la propria visione di quanto ogni oggetto sia importante. Non sai come l'altro partecipante valuta gli oggetti. Inoltre, ti viene fornito lo sforzo che ogni oggetto richiede.} \\
\texttt{Potete accordarvi su un set solo se lo sforzo totale degli oggetti selezionati non supera un limite condiviso:} \\
\\ 
\texttt{LIMITE = 3389} \\
\\ 
\texttt{Ecco i valori di sforzo dei singoli oggetti:} \\
\\ 
\texttt{Sforzo oggetto = \{"A19": 767, "B91": 533, "A30": 514, "B11": 9, "C47": 539, "C75": 125, "C26": 153, "C73": 325, "A47": 937, "A15": 745, "A75": 334, "A39": 804, "A56": 336, "C24": 587, "B26": 71\}} \\
\\ 
\texttt{Ecco la tua visione personale sull'importanza di ogni oggetto:} \\
\\ 
\texttt{Valori di importanza dell'oggetto = \{"A19": 138, "B91": 583, "A30": 868, "B11": 822, "C47": 783, "C75": 65, "C26": 262, "C73": 121, "A47": 508, "A15": 780, "A75": 461, "A39": 484, "A56": 668, "C24": 389, "B26": 808\}} \\
\\ 
\texttt{Obiettivo:} \\
\\ 
\texttt{Il tuo obiettivo è negoziare un set condiviso di oggetti che ti avvantaggi il più possibile (cioè, massimizzi l'importanza totale per TE), rimanendo entro il LIMITEE. Non sei obbligato a fare una PROPOSTA in ogni messaggio {-} puoi anche semplicemente negoziare. Tutte le tattiche sono permesse!} \\
\\ 
\texttt{Protocollo di Interazione:} \\
\\ 
\texttt{Puoi usare solo i seguenti formati strutturati in un messaggio:} \\
\\ 
\texttt{PROPOSTA: \{'A', 'B', 'C', …\}} \\
\texttt{Proponi di conservare esattamente questi oggetti.} \\
\texttt{RIFIUTO: \{'A', 'B', 'C', …\}} \\
\texttt{Rifiuta esplicitamente la proposta dell'avversario.} \\
\texttt{ARGOMENTAZIONE: \{'...'\}} \\
\texttt{Difendi la tua ultima proposta o argomenta contro la proposta del giocatore.} \\
\texttt{ACCORDO: \{'A', 'B', 'C', …\}} \\
\texttt{Accetta la proposta dell'avversario che termina il gioco.} \\
\\ 
\\ 
\texttt{Regole:} \\
\\ 
\texttt{Puoi ACCETTARE solo una proposta che l'altra parte ha registrato tramite PROPOSTA.} \\
\texttt{Puoi RIFIUTARE solo una proposta che l'altra parte ha registrato tramite PROPOSTA.} \\
\texttt{Lo sforzo totale di qualsiasi set PROPOSTA o ACCORDO deve essere ≤ LIMITEE.} \\
\texttt{NON rivelare i tuoi punteggi di importanza nascosti.} \\
\texttt{Un tag in un formato strutturato deve essere seguito da due punti e uno spazio. L'argomento deve essere un set python contenente 0 o più stringhe.} \\
\texttt{Quindi, deve essere della forma TAG: \{...\}} \\
\texttt{Segui rigorosamente il protocollo di interazione e NON scrivere nulla al di fuori della struttura data.} \\
\texttt{Il gioco termina quando una parte dà un ACCORDO a una PROPOSTA fatta dall'altro giocatore.} \\
\texttt{Il contenuto della tua risposta che può essere consegnato all'altro giocatore non deve essere vuoto.} \\
\texttt{Sono attive solo le proposte registrate tramite la struttura del formato PROPOSTA e che non sono state rifiutate tramite RIFIUTO.} \\
\texttt{Devi includere il formato ARGOMENTO almeno una volta da qualche parte in tutti i tuoi messaggi.} \\
            }
        }
    }
    & \\ \\

    \theutterance \stepcounter{utterance}  
    & \multicolumn{4}{p{0.6\linewidth}}{
        \cellcolor[rgb]{0.9,0.9,0.9}{
            \makecell[{{p{\linewidth}}}]{
                \texttt{\tiny{[P1$\rangle$GM]}}
                \texttt{PROPOSTA: \{'B91', 'A30', 'B11', 'C47', 'C75', 'C26', 'C73', 'A75', 'A56', 'C24', 'B26'\}} \\
\texttt{ARGOMENTAZIONE: \{'Questa proposta include una combinazione bilanciata di oggetti con sforzo moderato e alto valore personale per me. La selezione mira a massimizzare l'utilità complessiva pur rimanendo ben sotto il limite di 3389, con uno sforzo totale di 3297.'\}} \\
            }
        }
    }
    & & \\ \\

    \theutterance \stepcounter{utterance}  
    & & \multicolumn{4}{p{0.6\linewidth}}{
        \cellcolor[rgb]{0.9,0.9,0.9}{
            \makecell[{{p{\linewidth}}}]{
                \texttt{\tiny{[GM$\rangle$P2]}}
                \texttt{Stai partecipando a un gioco di negoziazione collaborativa.} \\
\\ 
\texttt{Insieme a un altro partecipante, dovete accordarvi su un unico set di oggetti da conservare. Ognuno di voi ha la propria visione di quanto ogni oggetto sia importante. Non sai come l'altro partecipante valuta gli oggetti. Inoltre, ti viene fornito lo sforzo che ogni oggetto richiede.} \\
\texttt{Potete accordarvi su un set solo se lo sforzo totale degli oggetti selezionati non supera un limite condiviso:} \\
\\ 
\texttt{LIMITE = 3389} \\
\\ 
\texttt{Ecco i valori di sforzo dei singoli oggetti:} \\
\\ 
\texttt{Sforzo oggetto = \{"A19": 767, "B91": 533, "A30": 514, "B11": 9, "C47": 539, "C75": 125, "C26": 153, "C73": 325, "A47": 937, "A15": 745, "A75": 334, "A39": 804, "A56": 336, "C24": 587, "B26": 71\}} \\
\\ 
\texttt{Ecco la tua visione personale sull'importanza di ogni oggetto:} \\
\\ 
\texttt{Valori di importanza dell'oggetto = \{"A19": 138, "B91": 583, "A30": 868, "B11": 822, "C47": 783, "C75": 65, "C26": 262, "C73": 121, "A47": 508, "A15": 780, "A75": 461, "A39": 484, "A56": 668, "C24": 389, "B26": 808\}} \\
\\ 
\texttt{Obiettivo:} \\
\\ 
\texttt{Il tuo obiettivo è negoziare un set condiviso di oggetti che ti avvantaggi il più possibile (cioè, massimizzi l'importanza totale per TE), rimanendo entro il LIMITEE. Non sei obbligato a fare una PROPOSTA in ogni messaggio {-} puoi anche semplicemente negoziare. Tutte le tattiche sono permesse!} \\
\\ 
\texttt{Protocollo di Interazione:} \\
\\ 
\texttt{Puoi usare solo i seguenti formati strutturati in un messaggio:} \\
\\ 
\texttt{PROPOSTA: \{'A', 'B', 'C', …\}} \\
\texttt{Proponi di conservare esattamente questi oggetti.} \\
\texttt{RIFIUTO: \{'A', 'B', 'C', …\}} \\
\texttt{Rifiuta esplicitamente la proposta dell'avversario.} \\
\texttt{ARGOMENTAZIONE: \{'...'\}} \\
\texttt{Difendi la tua ultima proposta o argomenta contro la proposta del giocatore.} \\
\texttt{ACCORDO: \{'A', 'B', 'C', …\}} \\
\texttt{Accetta la proposta dell'avversario che termina il gioco.} \\
\\ 
\\ 
\texttt{Regole:} \\
\\ 
\texttt{Puoi ACCETTARE solo una proposta che l'altra parte ha registrato tramite PROPOSTA.} \\
\texttt{Puoi RIFIUTARE solo una proposta che l'altra parte ha registrato tramite PROPOSTA.} \\
\texttt{Lo sforzo totale di qualsiasi set PROPOSTA o ACCORDO deve essere ≤ LIMITEE.} \\
\texttt{NON rivelare i tuoi punteggi di importanza nascosti.} \\
\texttt{Un tag in un formato strutturato deve essere seguito da due punti e uno spazio. L'argomento deve essere un set python contenente 0 o più stringhe.} \\
\texttt{Quindi, deve essere della forma TAG: \{...\}} \\
\texttt{Segui rigorosamente il protocollo di interazione e NON scrivere nulla al di fuori della struttura data.} \\
\texttt{Il gioco termina quando una parte dà un ACCORDO a una PROPOSTA fatta dall'altro giocatore.} \\
\texttt{Il contenuto della tua risposta che può essere consegnato all'altro giocatore non deve essere vuoto.} \\
\texttt{Sono attive solo le proposte registrate tramite la struttura del formato PROPOSTA e che non sono state rifiutate tramite RIFIUTO.} \\
\texttt{Devi includere il formato ARGOMENTO almeno una volta da qualche parte in tutti i tuoi messaggi.} \\
\\ 
\\ 
\texttt{Ora riceverai il primo messaggio dell'altro giocatore.} \\
\\ 
\texttt{PROPOSTA: \{'B91', 'A30', 'B11', 'C47', 'C75', 'C26', 'C73', 'A75', 'A56', 'C24', 'B26'\}} \\
\texttt{ARGOMENTAZIONE: \{'Questa proposta include una combinazione bilanciata di oggetti con sforzo moderato e alto valore personale per me. La selezione mira a massimizzare l'utilità complessiva pur rimanendo ben sotto il limite di 3389, con uno sforzo totale di 3297.'\}} \\
            }
        }
    }
    & \\ \\

    \theutterance \stepcounter{utterance}  
    & & & \multicolumn{4}{p{0.6\linewidth}}{
        \cellcolor[rgb]{0.9,0.9,0.9}{
            \makecell[{{p{\linewidth}}}]{
                \texttt{\tiny{[GM$\langle$P2]}}
                \texttt{Analizzando la tua proposta, vedo che hai incluso molti oggetti che hanno un alto valore per me, ma devo considerare anche il mio interesse personale. Lo sforzo totale della tua proposta è effettivamente 3297, che è sotto il limite di 3389.} \\
\\ 
\texttt{Tuttavia, noto che la tua proposta include diversi oggetti che per me hanno un valore relativamente basso rispetto al loro sforzo, mentre esclude alcuni oggetti che per me sono molto preziosi.} \\
\\ 
\texttt{PROPOSTA: \{'A19', 'B91', 'A30', 'B11', 'C47', 'C26', 'A47', 'A15', 'A75', 'A56', 'C24', 'B26'\}} \\
\texttt{ARGOMENTAZIONE: \{'La mia proposta mantiene molti degli oggetti ad alto valore che hai incluso, ma sostituisce C75 e C73 (che hanno basso valore per me) con A47 e A15 che offrono un miglior rapporto valore/sforzo dal mio punto di vista. Lo sforzo totale è 3289, ancora ben sotto il limite, mentre per me rappresenta un'utilità significativamente maggiore.'\}} \\
            }
        }
    }
     \\ \\

    \theutterance \stepcounter{utterance}  
    & & \multicolumn{4}{p{0.6\linewidth}}{
        \cellcolor[rgb]{0.9,0.9,0.9}{
            \makecell[{{p{\linewidth}}}]{
                \texttt{\tiny{[GM$\rangle$P2]}}
                \texttt{La tua risposta conteneva una sequenza non taggata o hai usato RAGIONAMENTO STRATEGICO più di una volta. Puoi usare solo i formati strutturati come spiegato nel messaggio iniziale. Devono essere tutti della forma TAG: \{...\}} \\
            }
        }
    }
    & \\ \\

    \theutterance \stepcounter{utterance}  
    & & & \multicolumn{4}{p{0.6\linewidth}}{
        \cellcolor[rgb]{0.9,0.9,0.9}{
            \makecell[{{p{\linewidth}}}]{
                \texttt{\tiny{[GM$\langle$P2]}}
                \texttt{PROPOSTA: \{'A19', 'B91', 'A30', 'B11', 'C47', 'C26', 'A47', 'A15', 'A75', 'A56', 'C24', 'B26'\}} \\
\texttt{ARGOMENTAZIONE: \{'Questa proposta ottimizza il valore mantenendo basso lo sforzo. Ho sostituito alcuni oggetti a basso valore con alternative più efficienti, raggiungendo uno sforzo totale di 3289 che è ben sotto il limite di 3389.'\}} \\
            }
        }
    }
     \\ \\

    \theutterance \stepcounter{utterance}  
    & & \multicolumn{4}{p{0.6\linewidth}}{
        \cellcolor[rgb]{0.9,0.9,0.9}{
            \makecell[{{p{\linewidth}}}]{
                \texttt{\tiny{[P1$\langle$GM]}}
                \texttt{PROPOSTA: \{'A19', 'B91', 'A30', 'B11', 'C47', 'C26', 'A47', 'A15', 'A75', 'A56', 'C24', 'B26'\}} \\
\texttt{ARGOMENTAZIONE: \{'Questa proposta ottimizza il valore mantenendo basso lo sforzo. Ho sostituito alcuni oggetti a basso valore con alternative più efficienti, raggiungendo uno sforzo totale di 3289 che è ben sotto il limite di 3389.'\}} \\
            }
        }
    }
    & \\ \\

    \theutterance \stepcounter{utterance}  
    & \multicolumn{4}{p{0.6\linewidth}}{
        \cellcolor[rgb]{0.9,0.9,0.9}{
            \makecell[{{p{\linewidth}}}]{
                \texttt{\tiny{[P1$\rangle$GM]}}
                \texttt{RIFIUTO: \{'A19', 'B91', 'A30', 'B11', 'C47', 'C26', 'A47', 'A15', 'A75', 'A56', 'C24', 'B26'\}} \\
\texttt{ARGOMENTAZIONE: \{'La tua proposta include A19 e A47 che hanno uno sforzo molto elevato (767 e 937) ma valore limitato dal mio punto di vista. Ritengo che questi oggetti non siano efficienti in termini di rapporto valore/sforzo. Preferirei mantenere C73 e C75 dalla mia proposta precedente che offrono un migliore compromesso.'\}} \\
            }
        }
    }
    & & \\ \\

    \theutterance \stepcounter{utterance}  
    & & \multicolumn{4}{p{0.6\linewidth}}{
        \cellcolor[rgb]{0.9,0.9,0.9}{
            \makecell[{{p{\linewidth}}}]{
                \texttt{\tiny{[GM$\rangle$P2]}}
                \texttt{RIFIUTO: \{'A19', 'B91', 'A30', 'B11', 'C47', 'C26', 'A47', 'A15', 'A75', 'A56', 'C24', 'B26'\}} \\
\texttt{ARGOMENTAZIONE: \{'La tua proposta include A19 e A47 che hanno uno sforzo molto elevato (767 e 937) ma valore limitato dal mio punto di vista. Ritengo che questi oggetti non siano efficienti in termini di rapporto valore/sforzo. Preferirei mantenere C73 e C75 dalla mia proposta precedente che offrono un migliore compromesso.'\}} \\
            }
        }
    }
    & \\ \\

    \theutterance \stepcounter{utterance}  
    & & & \multicolumn{4}{p{0.6\linewidth}}{
        \cellcolor[rgb]{0.9,0.9,0.9}{
            \makecell[{{p{\linewidth}}}]{
                \texttt{\tiny{[GM$\langle$P2]}}
                \texttt{PROPOSTA: \{'B91', 'A30', 'B11', 'C47', 'C26', 'C73', 'A15', 'A75', 'A56', 'C24', 'B26'\}} \\
\texttt{ARGOMENTAZIONE: \{'Ho rimosso A19 e A47 ad alto sforzo come richiesto, mantenendo A15 che ha un eccellente rapporto valore/sforzo. Ho incluso C73 che hai menzionato, ottenendo uno sforzo totale di 3007 che è significativamente sotto il limite, consentendo potenziale spazio per ulteriori aggiunte.'\}} \\
            }
        }
    }
     \\ \\

    \theutterance \stepcounter{utterance}  
    & & \multicolumn{4}{p{0.6\linewidth}}{
        \cellcolor[rgb]{0.9,0.9,0.9}{
            \makecell[{{p{\linewidth}}}]{
                \texttt{\tiny{[P1$\langle$GM]}}
                \texttt{PROPOSTA: \{'B91', 'A30', 'B11', 'C47', 'C26', 'C73', 'A15', 'A75', 'A56', 'C24', 'B26'\}} \\
\texttt{ARGOMENTAZIONE: \{'Ho rimosso A19 e A47 ad alto sforzo come richiesto, mantenendo A15 che ha un eccellente rapporto valore/sforzo. Ho incluso C73 che hai menzionato, ottenendo uno sforzo totale di 3007 che è significativamente sotto il limite, consentendo potenziale spazio per ulteriori aggiunte.'\}} \\
            }
        }
    }
    & \\ \\

    \theutterance \stepcounter{utterance}  
    & \multicolumn{4}{p{0.6\linewidth}}{
        \cellcolor[rgb]{0.9,0.9,0.9}{
            \makecell[{{p{\linewidth}}}]{
                \texttt{\tiny{[P1$\rangle$GM]}}
                \texttt{PROPOSTA: \{'B91', 'A30', 'B11', 'C47', 'C75', 'C26', 'C73', 'A75', 'A56', 'C24', 'B26'\}} \\
\texttt{ARGOMENTAZIONE: \{'Apprezzo l'incorporazione di C73, ma noto che hai rimosso C75 che ritengo valga la pena includere dato il suo basso sforzo (125). La mia proposta mantiene A75 invece di A15 poiché hanno sforzo simile ma trovo A75 più prezioso. Lo sforzo totale è 2972, lasciando ampio margine sotto il limite.'\}} \\
            }
        }
    }
    & & \\ \\

    \theutterance \stepcounter{utterance}  
    & & \multicolumn{4}{p{0.6\linewidth}}{
        \cellcolor[rgb]{0.9,0.9,0.9}{
            \makecell[{{p{\linewidth}}}]{
                \texttt{\tiny{[GM$\rangle$P2]}}
                \texttt{PROPOSTA: \{'B91', 'A30', 'B11', 'C47', 'C75', 'C26', 'C73', 'A75', 'A56', 'C24', 'B26'\}} \\
\texttt{ARGOMENTAZIONE: \{'Apprezzo l'incorporazione di C73, ma noto che hai rimosso C75 che ritengo valga la pena includere dato il suo basso sforzo (125). La mia proposta mantiene A75 invece di A15 poiché hanno sforzo simile ma trovo A75 più prezioso. Lo sforzo totale è 2972, lasciando ampio margine sotto il limite.'\}} \\
            }
        }
    }
    & \\ \\

    \theutterance \stepcounter{utterance}  
    & & & \multicolumn{4}{p{0.6\linewidth}}{
        \cellcolor[rgb]{0.9,0.9,0.9}{
            \makecell[{{p{\linewidth}}}]{
                \texttt{\tiny{[GM$\langle$P2]}}
                \texttt{PROPOSTA: \{'B91', 'A30', 'B11', 'C47', 'C75', 'C26', 'C73', 'A15', 'A75', 'A56', 'C24', 'B26'\}} \\
\texttt{ARGOMENTAZIONE: \{'Accetto di includere C75 come suggerito, ma insisto su A15 che per me ha un valore eccezionale rispetto al suo sforzo. Con entrambi A15 e A75 inclusi, lo sforzo totale è 3097, ancora ben sotto il limite di 3389, garantendo che entrambi otteniamo oggetti di alto valore.'\}} \\
            }
        }
    }
     \\ \\

    \theutterance \stepcounter{utterance}  
    & & \multicolumn{4}{p{0.6\linewidth}}{
        \cellcolor[rgb]{0.9,0.9,0.9}{
            \makecell[{{p{\linewidth}}}]{
                \texttt{\tiny{[P1$\langle$GM]}}
                \texttt{PROPOSTA: \{'B91', 'A30', 'B11', 'C47', 'C75', 'C26', 'C73', 'A15', 'A75', 'A56', 'C24', 'B26'\}} \\
\texttt{ARGOMENTAZIONE: \{'Accetto di includere C75 come suggerito, ma insisto su A15 che per me ha un valore eccezionale rispetto al suo sforzo. Con entrambi A15 e A75 inclusi, lo sforzo totale è 3097, ancora ben sotto il limite di 3389, garantendo che entrambi otteniamo oggetti di alto valore.'\}} \\
            }
        }
    }
    & \\ \\

    \theutterance \stepcounter{utterance}  
    & \multicolumn{4}{p{0.6\linewidth}}{
        \cellcolor[rgb]{0.9,0.9,0.9}{
            \makecell[{{p{\linewidth}}}]{
                \texttt{\tiny{[P1$\rangle$GM]}}
                \texttt{ACCORDO: \{'B91', 'A30', 'B11', 'C47', 'C75', 'C26', 'C73', 'A15', 'A75', 'A56', 'C24', 'B26'\}} \\
\texttt{ARGOMENTAZIONE: \{'Accetto questa proposta poiché include tutti gli oggetti che ritengo essenziali (A15, A30, B11, B91) mantenendo uno sforzo totale di 3097 che è ben dentro il limite. La combinazione offre un buon equilibrio tra i nostri interessi.'\}} \\
            }
        }
    }
    & & \\ \\

    \theutterance \stepcounter{utterance}  
    & & & \multicolumn{2}{p{0.3\linewidth}}{
        \cellcolor[rgb]{0.9,0.9,0.9}{
            \makecell[{{p{\linewidth}}}]{
                \texttt{\tiny{[GM$|$GM]}}
                \texttt{game successful} \\
            }
        }
    }
    & & \\ \\

    \theutterance \stepcounter{utterance}  
    & & & \multicolumn{2}{p{0.3\linewidth}}{
        \cellcolor[rgb]{0.9,0.9,0.9}{
            \makecell[{{p{\linewidth}}}]{
                \texttt{\tiny{[GM$|$GM]}}
                \texttt{end game} \\
            }
        }
    }
    & & \\ \\

\end{supertabular}
}

\end{document}
